% ==============================================================================
% TRABALHOS RELACIONADOS
% Fontes verificadas nos artigos:
%   ramirez2017sudden — PLoS ONE 12(10):e0186152 (texto completo)
%   kim2025deep       — J Clin Med 14(20):7209 (texto completo)
%   swee2026multimodal — números já usados em resultados_marcadores.tex;
%                        o preprint não pôde ser reaberto agora.
%                        [conferir] = completar com o PDF do preprint.
% ==============================================================================

\chapter{Trabalhos Relacionados}
\label{cap:trabalhos-relacionados}

A predição de morte súbita cardíaca a partir do \gls{ECG} ambulatorial tem
seguido dois caminhos. O primeiro, mais antigo, extrai do registro de Holter
marcadores definidos por especialistas --- variabilidade da frequência
cardíaca, turbulência da frequência, alternância da onda T, dinâmica da
repolarização, contagem de arritmias --- e os combina com variáveis clínicas
em modelos estatísticos. O segundo, mais recente, aplica redes neurais
profundas diretamente ao sinal, na expectativa de que a rede encontre padrões
que os marcadores clássicos não capturam. Esse segundo caminho se apoia no
sucesso das redes convolucionais em tarefas diagnósticas sobre o \gls{ECG},
como a detecção de arritmias no Holter \cite{hannun2019cardiologist} e o
diagnóstico automático no \gls{ECG} de 12 derivações
\cite{ribeiro2020automatic}. Há, contudo, uma diferença importante entre
essas tarefas e a predição de morte súbita: nas tarefas diagnósticas, o
padrão a ser reconhecido está presente no traçado e há dezenas de milhares
de exemplos rotulados; na predição de um evento futuro, o padrão pode não
existir no traçado e os eventos são poucos.

Este capítulo apresenta três trabalhos que representam esses caminhos e que
servem de referência para a interpretação dos resultados desta dissertação.
O primeiro e o segundo utilizam a mesma base de dados deste trabalho, a
coorte MUSIC \cite{martinyebra2024music}: um com marcadores clássicos e
modelos de Cox, o outro com marcadores de segmentos curtos e aprendizado de
máquina. O terceiro aplica aprendizado profundo ao Holter de 24~horas em uma
coorte coreana de pacientes com infarto ou insuficiência cardíaca. O
capítulo se encerra com uma comparação entre eles e com o posicionamento
desta dissertação.

\section{Marcadores eletrocardiográficos e clínicos na coorte MUSIC}
\label{sec:trab-ramirez}

\citeonline{ramirez2017sudden} construíram modelos de risco para morte
súbita e para morte por falência de bomba na coorte MUSIC, combinando
marcadores do Holter e variáveis clínicas. A análise incluiu os pacientes em
ritmo sinusal com todos os índices calculáveis --- 597 dos 992 pacientes
originais ---, seguidos por mediana de 44~meses, período em que ocorreram
49 mortes súbitas e 62 mortes por falência de bomba. Os pacientes com
fibrilação atrial, ritmo de marca-passo ou ausência de extrassístoles
ventriculares foram excluídos, este último critério porque a turbulência da
frequência depende de extrassístoles para ser calculada.

Os marcadores eletrocardiográficos incluíram a turbulência da frequência (HRT)\glsadd{HRT}\glsunset{HRT}, a alternância da onda T (TWA)\glsadd{TWA}\glsunset{TWA}, índices de restituição do intervalo QT e do intervalo
entre o pico e o fim da onda T, a restituição da morfologia da onda T, a
amplitude dos intervalos RR e um índice de arritmia complexa. Para cada
desfecho, os autores ajustaram três modelos de Cox --- clínico, apenas
eletrocardiográfico e combinado ---, com seleção das variáveis por
eliminação regressiva, e derivaram deles escores de risco em pontos. O
modelo combinado para morte súbita reteve sexo masculino, classe III da
\gls{NYHA}, fração de ejeção $\leq$ 35\,\% e três marcadores da repolarização.
Os pacientes no quintil de maior risco tiveram risco de morte súbita 13,8
vezes maior que os dos dois quintis de menor risco, contra 4,0 vezes com o
modelo apenas clínico. No subgrupo com fração de ejeção reduzida, a
\gls{AUC} para morte súbita foi de 0,73 [0,64--0,81] com o modelo
combinado.

Os próprios autores apontam a principal limitação do trabalho: não houve
coorte de validação separada, pelo número reduzido de eventos. A seleção
das variáveis e a estimativa do desempenho foram feitas sobre os mesmos
pacientes, o que tende a superestimar a discriminação. Ainda assim, o
trabalho estabelece duas referências para esta dissertação: a restrição ao
ritmo sinusal como critério da coorte e a evidência de que, na MUSIC, a
informação prognóstica do Holter para morte súbita está concentrada em
marcadores da repolarização, que exigem delineação precisa da onda T.

\section{Integração de marcadores de segmentos curtos e variáveis clínicas}
\label{sec:trab-swee}

\citeonline{swee2026multimodal} também utilizaram a coorte MUSIC, com o
objetivo de prever morte súbita e morte por falência de bomba a partir da
integração de marcadores do Holter com variáveis clínicas. Em vez de
marcadores calculados sobre o registro inteiro, os autores extraíram
marcadores de segmentos de 30~segundos não sobrepostos, com detecção dos
picos R na derivação Y, e treinaram modelos de aprendizado de máquina sobre
essas representações. O desenho desses segmentos foi adotado no
protocolo~B desta dissertação (\autoref{subsec:met-protocolo-b}).

O trabalho incluiu 746 pacientes, sem restrição ao ritmo sinusal, e relata
\gls{AUC} de 0,682 $\pm$ 0,030 para o modelo apenas com marcadores do
\gls{ECG}, expressa como média entre as partições da validação cruzada.
% [conferir no preprint: número de eventos de morte súbita, modelos
% empregados, AUC do modelo clínico e do modelo combinado.]
Dois aspectos limitam a comparação com os resultados desta dissertação. O
primeiro é a coorte: sem a restrição ao ritmo sinusal, ela inclui pacientes
em fibrilação atrial, ritmo facilmente reconhecível no sinal e associado a
pior prognóstico, o que pode por si só elevar a discriminação, como mostrou
a análise de sensibilidade da \autoref{subsec:res-redes-ritmo}. O segundo é a
métrica: a média entre partições não equivale à \gls{AUC} das predições
agrupadas usada como medida principal neste trabalho.

\section{Aprendizado profundo no Holter de 24 horas}
\label{sec:trab-kim}

\citeonline{kim2025deep} aplicaram aprendizado profundo ao Holter de
24~horas de 1\,108 pacientes hospitalizados por infarto agudo do miocárdio
(885) ou insuficiência cardíaca aguda (223), provenientes de um registro
prospectivo em 25 centros da Coreia do Sul. O desfecho foi composto por
morte cardíaca e arritmias ventriculares em três anos, com 70 eventos.

O modelo não foi treinado de ponta a ponta sobre o sinal do Holter. Os
autores partiram de um codificador do tipo \textit{Transformer}
pré-treinado como autocodificador mascarado sobre \gls{ECG}s de
12~derivações de 10~segundos e, em seguida, ajustado para estimar a fração
de ejeção, usada como rótulo substituto. Com o codificador congelado, cada
registro de 24~horas da derivação V1, amostrado a 125~Hz, foi dividido em
9\,600 segmentos de 9~segundos, cada um representado por um vetor latente.
Somente a etapa final --- uma ResNet unidimensional compacta sobre a
sequência de vetores --- foi treinada com o desfecho, em validação cruzada
de cinco partições por paciente.

Na coorte completa, o modelo obteve \gls{AUC} de 0,74 [0,70--0,77] para o
desfecho composto e 0,79 [0,74--0,84] para morte cardíaca, superando a
turbulência da frequência (0,62) e a alternância da onda T (0,55), e
comparável à fração de ejeção isolada (0,75). No subgrupo de insuficiência
cardíaca, no entanto, a \gls{AUC} caiu para 0,60 [0,53--0,68], igual à da
turbulência da frequência, e a combinação com a fração de ejeção não
acrescentou desempenho (0,58). Os autores reconhecem o número reduzido de
eventos e a ausência de validação externa.

Dois elementos desse trabalho são particularmente relevantes aqui. O
primeiro é a estratégia: o extrator de características foi aprendido em
grandes bases de \gls{ECG} de repouso e com um rótulo substituto
abundante, e apenas um classificador pequeno foi ajustado com os poucos
eventos disponíveis --- o oposto do treinamento de ponta a ponta adotado
nas redes desta dissertação. O segundo é o resultado no subgrupo de
insuficiência cardíaca, população equivalente à da coorte MUSIC, em que
mesmo esse modelo pré-treinado ficou próximo do limite inferior de
discriminação útil.

\section{Comparação e posicionamento deste trabalho}
\label{sec:trab-comparacao}

O \autoref{qua:trab-comparacao} resume os três trabalhos e esta dissertação.

\begin{quadro}[h!]
	\centering
	\Caption{\label{qua:trab-comparacao} Comparação entre os trabalhos
		relacionados e esta dissertação}
	\UFCqua{}{
		\footnotesize
		\begin{tabular}{|p{2.3cm}|p{2.6cm}|p{2.9cm}|p{2.9cm}|p{3.1cm}|}
			\hline
			\textbf{Trabalho}                                                                                & \textbf{Coorte}              & \textbf{Entrada} &
			\textbf{Método e validação}                                                                      & \textbf{Resultado principal}                      \\
			\hline
			\citeonline{ramirez2017sudden}                                                                   &
			MUSIC, ritmo sinusal; 597 pacientes, 49 mortes súbitas                                           &
			Marcadores do Holter (repolarização, turbulência) e clínicos                                     &
			Cox com seleção regressiva; sem validação separada                                               &
			Razão de risco 13,8 (alto contra baixo risco); \gls{AUC} 0,73 na fração de ejeção reduzida                                                           \\
			\hline
			\citeonline{swee2026multimodal}                                                                  &
			MUSIC, sem restrição de ritmo; 746 pacientes                                                     &
			Marcadores de segmentos de 30~s e clínicos                                                       &
			Aprendizado de máquina; validação cruzada                                                        &
			\gls{AUC} 0,682 (apenas \gls{ECG}, média entre partições)                                                                                            \\
			\hline
			\citeonline{kim2025deep}                                                                         &
			Coreia; infarto ou insuficiência cardíaca; 1\,108 pacientes, 70 eventos                          &
			Holter de 24~h, derivação V1                                                                     &
			Codificador pré-treinado e congelado; validação cruzada em 5 partições                           &
			\gls{AUC} 0,74 na coorte; 0,60 na insuficiência cardíaca                                                                                             \\
			\hline
			Esta dissertação                                                                                 &
			MUSIC, ritmo sinusal; 517 pacientes, 55 mortes súbitas em 4 anos                                 &
			Sinal das três derivações; marcadores do Holter; clínicos                                        &
			Redes treinadas de ponta a ponta; modelos lineares e florestas; validação cruzada em 5 partições &
			\gls{AUC} 0,48--0,56 (redes); 0,661 (clínicos)                                                                                                       \\
			\hline
		\end{tabular}
	}{
		\Fonte{elaborado pelo autor.}
	}
\end{quadro}

A comparação revela um padrão. Os desempenhos mais altos foram obtidos em
condições que favorecem a discriminação: avaliação nos mesmos dados usados
para selecionar as variáveis \cite{ramirez2017sudden}, coortes que incluem
ritmos não sinusais \cite{swee2026multimodal} ou populações em que
predominam pacientes pós-infarto, nas quais a fração de ejeção já
discrimina bem \cite{kim2025deep}. Quando a população se restringe a
pacientes com insuficiência cardíaca, como no subgrupo de
\citeonline{kim2025deep}, o desempenho se aproxima de 0,60.

Esta dissertação se diferencia dos trabalhos anteriores em três aspectos. O
primeiro é a avaliação, no mesmo protocolo e sobre os mesmos pacientes, de
três abordagens --- redes treinadas sobre o sinal, marcadores extraídos do
Holter e variáveis clínicas ---, o que permite compará-las diretamente. O
segundo é a ênfase no controle do otimismo: validação cruzada por paciente,
seleção de variáveis dentro de cada partição, \gls{AUC} das predições
agrupadas com intervalos de confiança e o registro das configurações que
não funcionaram, em vez de apenas da melhor
\cite{kapoor2023leakage}. O terceiro é a restrição ao ritmo sinusal, que
elimina o atalho oferecido pela fibrilação atrial e torna a tarefa mais
próxima do problema clínico de interesse: identificar, entre pacientes com
insuficiência cardíaca em ritmo sinusal, aqueles com maior risco de morte
súbita.

% \chapter{Trabalhos Relacionados}
% \label{cap:trabalhos-relacionados}

% Nesta seção, escreva sobre uma introdução sobre os trabalhos relacionados se eles existirem. Evite começar da seção secundária, ou seja, escreva um texto para introduzir as subseções subsequentes, as quais serão sobre cada trabalho relacionado.

% \section{Trabalho Relacionado A}
% \label{sec:trabalho-relacionado-a}

% Insira o texto sobre o Trabalho A aqui:

% Texto texto texto texto texto texto texto texto texto texto texto texto texto texto texto texto texto texto texto.

% \section{Trabalho Relacionado B}
% \label{sec:trabalho-relacionado-b}

% Insira o texto sobre o Trabalho B aqui:

% Texto texto texto texto texto texto texto texto texto texto texto texto texto texto texto texto texto texto.

% Exemplos de inserção de quadro:


% 	\begin{quadro}[h!]	
% 		\centering
% 		\Caption{\label{qua:exemplo-1} Praesent ex velit, pulvinar at massa vel, fermentum dictum mauris. Ut feugiat accumsan augue}		
% 		\UFCqua{}{
% 			\begin{tabular}{|c|c|l|l|}
% 				\hline
% 				Quisque & pharetra & tempus & vulputate \\
% 				\hline
% 				E1 & Complete coverage by a single transcript & Both  & Complete\\
% 				\hline
% 				E2 & Complete coverage by more than & Both splice sites & Complete\\
% 				\hline
% 				E3 & Partial coverage & Both splice sites & Both \\				
% 				\hline
% 			\end{tabular}
% 		}{
% 			\Fonte{elaborado pelo autor.}
% 		}
% 	\end{quadro}



% 	\begin{quadro}[h!]	
% 		\centering
% 		\Caption{\label{qua:exemplo-2} Duis faucibus, enim quis tincidunt pellentesque}		
% 		\UFCqua{}{
% 			\begin{tabular}{|c|c|}
% 				\hline
% 				Quisque & pharetra \\
% 				\hline
% 				E1 & Complete coverage by a single transcript \\
% 				\hline
% 				E2 & Complete coverage by more than \\
% 				\hline
% 				E3 & Partial coverage \\
% 				\hline
% 				E4 & Partial coverage \\
% 				\hline
% 				E5 & Partial coverage \\
% 				\hline
% 				E6 & Partial coverage \\
% 				\hline
% 				E7 & Partial coverage \\
% 				\hline
% 			\end{tabular}
% 		}{
% 			\Fonte{elaborado pelo autor.}
% 		}
% 	\end{quadro}


% Texto  texto texto texto texto.

% %\Gls{ambiguidade}
% %\Gls{braile}
% %\Gls{coerencia}
% %\Gls{dialetos}
% %\Gls{elipse}
% %\Gls{locucao-adjetiva}
% %\Gls{modificadores}
% %\Gls{paronimos}
% %\Gls{sintese}
% %\Gls{borboleta}